\chapter{正规族}
\label{chap:normal-families}

逐点有界并不能保证一列全纯函数有收敛子列。复分析中起作用的是局部一致
有界：Cauchy 估计把函数值的一致界传给导数，于是一个函数族在每个紧集上
同时具有有界性和等度连续性。Arzel\`a--Ascoli 定理随后给出收敛子列，而
Cauchy 积分公式保证极限仍全纯。这就是 Montel 定理的内容。

正规族方法不直接计算候选极限。它先从任意序列中抽取极限，再用归一化、
恒等定理或极值条件识别极限。下一章的 Riemann 映射定理将采用这一方法。

\section{局部一致收敛和全纯极限}

设 $\Omega\subset\C$ 为区域。称 $f_n\to f$ 在 $\Omega$ 上
\term{局部一致收敛}，如果对每个紧集 $K\Subset\Omega$，都有
\[
 \sup_{z\in K}|f_n(z)-f(z)|\longrightarrow0.
\]
这也称为紧开收敛。它比整个 $\Omega$ 上的一致收敛弱，但足以在任意固定
闭曲线和任意内含圆盘上使用一致极限。

\begin{theorem}[Weierstrass 全纯极限定理]\label{thm:holomorphic-local-uniform-limit}
若 $f_n$ 在 $\Omega$ 上全纯且局部一致收敛到 $f$，则 $f$ 全纯。并且对每个
$k\ge0$，
\[
 f_n^{(k)}\longrightarrow f^{(k)}
\]
也在 $\Omega$ 上局部一致。
\end{theorem}

\begin{proof}
固定 $a\in\Omega$，取 $R>0$ 使闭圆盘
$\overline{B(a,R)}\subset\Omega$。在较小圆盘 $|z-a|<R$ 中，Cauchy 积分
公式给
\[
 f_n(z)=\frac{1}{2\pi i}\int_{|\zeta-a|=R}
 \frac{f_n(\zeta)}{\zeta-z}\,d\zeta.
\]
$f_n$ 在圆周上一致收敛到 $f$，故可在积分中取极限，得到同一公式，其中
$f_n$ 换成 $f$。右侧关于 $z$ 全纯，所以 $f$ 在 $B(a,R)$ 中全纯。

再对核求导，得到
\[
 f_n^{(k)}(z)-f^{(k)}(z)
 =\frac{k!}{2\pi i}\int_{|\zeta-a|=R}
 \frac{f_n(\zeta)-f(\zeta)}{(\zeta-z)^{k+1}}\,d\zeta.
\]
若 $|z-a|\le r<R$，右侧绝对值至多为
\[
 \frac{k!R}{(R-r)^{k+1}}
 \max_{|\zeta-a|=R}|f_n(\zeta)-f(\zeta)|,
\]
因而在每个较小闭圆盘上一致趋零。任意紧集可由有限个这样的圆盘覆盖，结论
成立。
\end{proof}

定理中的局部一致性不能在证明中直接换成逐点收敛：逐点收敛既不能保证在
Cauchy 积分中交换极限，也不给导数的统一控制。若再有局部有界性，并且已知
函数值在一个具有内部聚点的集合上收敛，Vitali--Porter 定理才会把这些较弱
信息提升为局部一致收敛。

在单位圆盘中，$f_n(z)=z^n$ 局部一致收敛到零，因为对
$K\Subset\mathbb D$ 存在 $r<1$ 使 $|z|\le r$ 于 $K$，从而
$\sup_K|z|^n\le r^n$。它在闭单位圆盘上并不一致收敛。这一例子也说明
“局部”不能省略：紧集必须与区域边界保持正距离。

\section{正规族}

\begin{definition}
区域 $\Omega$ 上的一族全纯函数 $\mathcal F$ 称为\term{正规族}，如果
$\mathcal F$ 中任意序列都含有一个在 $\Omega$ 上局部一致收敛的子列，且
极限为有限值全纯函数。
\end{definition}

这里采用有限值版本；常数列 $f_n\equiv n$ 因而没有正规子列。有些文献在
Riemann 球面的球面距离下允许子列收敛到无穷远点，那是亚纯函数族常用的
另一种约定，使用时必须另行说明。

两个简单的函数族可以说明定义中各项的含义。族
\[
 \mathcal F_1=\{az:|a|<1\}
\]
在 $\mathbb D$ 上正规；系数列 $(a_n)$ 总可抽取收敛子列，函数列便在每个紧集
上一致收敛。它却不是紧开拓扑下的闭集，因为 $a_n\to1$ 时极限 $z$ 不属于
$\mathcal F_1$。所以正规性说的是相对紧性，并不包含封闭性。相反，
$\mathcal F_2=\{nz:n\ge1\}$ 在原点逐点有界，却在任何含有非零点的紧集上
都不一致有界，因而不是正规族。正规性所需的控制必须同时作用于一个邻域，
只控制单个点没有作用。

取紧集穷竭
\[
 K_m=\{z\in\Omega:|z|\le m,
 \operatorname{dist}(z,\C\setminus\Omega)\ge1/m\},
\]
这里约定到空集的距离为无穷，所以也包括 $\Omega=\C$。
这些集合闭且有界，且 $K_m\subset\operatorname{int}K_{m+1}$：
两条界限从 $m$ 到 $m+1$ 都严格放宽。每点到补集的距离为正，
因而各 $K_m$ 的内部覆盖 $\Omega$。删去开头的空项并重新编号后，
得到非空紧集穷竭。任意紧集由这些内部的有限多个覆盖，故包含于某个 $K_m$。
在全纯函数空间上定义
\[
 d(f,g)=\sum_{m=1}^{\infty}2^{-m}
 \min\{1,\|f-g\|_{C(K_m)}\}.
\]
级数各项不超过 $2^{-m}$，所以 $d$ 有限。对称性和非负性直接成立；
若 $d(f,g)=0$，则 $f=g$ 于每个 $K_m$，从而在 $\Omega$ 上相等。
由上确界范数的三角不等式及
$\min(1,s+t)\le\min(1,s)+\min(1,t)$，逐项相加得到 $d$ 的三角不等式。
因此它确为度量。先控制有限项、再控制级数尾部可知，
$d(f_n,f)\to0$ 当且仅当 $f_n\to f$ 局部一致。这个度量把“每个紧集
上的一致收敛”变成普通的序列收敛，也使逐层抽取的对角线过程可以同时控制
所有紧集。

称 $\mathcal F$ \term{局部有界}，如果每个点 $a\in\Omega$ 都有邻域 $U$ 和
常数 $M_U$，使
\[
 |f(z)|\le M_U\qquad(f\in\mathcal F, z\in U).
\]
有限覆盖说明，这等价于对每个 $K\Subset\Omega$ 都有
$\sup_{f\in\mathcal F}\sup_K|f|<\infty$。

\section{Montel 定理}

\begin{theorem}[Montel 定理]\label{thm:montel-locally-bounded}
区域上的任意局部有界全纯函数族都是正规族。
\end{theorem}

\begin{proof}
先在一个紧集上取得等度连续性。固定 $K\Subset\Omega$，选取开集 $U$ 使
$K\Subset U\Subset\Omega$。局部有界性与 $\overline U$ 的紧性给常数 $M$，
使 $|f|\le M$ 对所有 $f\in\mathcal F$ 及 $z\in U$ 成立。记
$d=\operatorname{dist}(K,\partial U)>0$，并取
\[
 0<\rho<\frac d2.
\]
令 $K_\rho=\{z:\operatorname{dist}(z,K)\le\rho\}$。对
$z\in K_\rho$，闭圆盘 $\overline{B(z,\rho)}$ 仍包含在 $U$ 中，故 Cauchy
导数估计给
\[
 |f'(z)|\le \frac{M}{\rho}.
\]
若 $z,w\in K$ 且 $|z-w|<\rho$，连接它们的线段包含在 $K_\rho$ 中，沿线段
积分便得
\[
 |f(z)-f(w)|\le \frac{M}{\rho}|z-w|.
\]
这就证明了 $\mathcal F|_K$ 等度连续；它还一致有界。

现取上一节的紧集穷竭 $K_1\subset K_2\subset\cdots$，以及任意序列
$(f_n)\subset\mathcal F$。由 Arzel\`a--Ascoli 定理，在 $K_1$ 上可抽取
一致收敛子列；从该子列中再抽取在 $K_2$ 上一致收敛的子列，如此继续。
取第 $m$ 层子列的第 $m$ 项组成对角子列 $(f_{n_m})$。对每个固定 $j$，它从
某项起属于第 $j$ 层子列，故在 $K_j$ 上一致收敛。各层极限在交叠处相同，
于是得到 $\Omega$ 上的局部一致极限 $f$。由
定理\ref{thm:holomorphic-local-uniform-limit}，$f$ 全纯。
\end{proof}

在这里，Cauchy 估计承担了复分析特有的一步。对任意连续函数族，局部一致
有界并不能推出等度连续；全纯性把函数值的控制自动转化为导数控制。
证明的其余两步各有不同作用：Arzel\`a--Ascoli 只在一个固定紧集上给出子列，
对角线过程把可数个紧集上的选择合并为同一子列；最后的 Weierstrass 定理才
保证所得连续极限仍属全纯函数类。缺少其中任何一步，都还没有得到区域上的
正规性。

有限值意义下，Montel 定理还有逆向结论。

\begin{proposition}\label{prop:normal-iff-locally-bounded}
一族全纯函数正规，当且仅当它局部有界。
\end{proposition}

\begin{proof}
只需证明正规蕴含局部有界。若在某个紧集 $K\Subset\Omega$ 上不一致有界，
可选 $f_n\in\mathcal F$ 和 $z_n\in K$ 使 $|f_n(z_n)|\ge n$。取子列使
$z_n\to z\in K$。正规性给进一步子列局部一致收敛到有限全纯函数，因而在
$K$ 的一个紧邻域上一致有界，与 $|f_n(z_n)|\to\infty$ 矛盾。
\end{proof}

特别地，所有从 $\Omega$ 映入单位圆盘的全纯函数构成正规族。若再施加在一点
取值、导数或单叶性等条件，需要检查这些条件在局部一致极限下是否封闭。
取值与导数条件由
定理\ref{thm:holomorphic-local-uniform-limit}保持；单叶性则由
推论\ref{cor:univalent-limit}保持到“单叶或常值”这一准确形式。

\section{Vitali--Porter 收敛定理}

Montel 定理只给子列。若函数列已在一个具有内部聚点的集合上收敛，恒等定理
会迫使所有子列极限相同，于是整个序列收敛。

\begin{theorem}[Vitali--Porter 定理]\label{thm:vitali-porter}
设 $(f_n)$ 是区域 $\Omega$ 上局部有界的全纯函数列。若存在集合
$E\subset\Omega$，$E$ 在 $\Omega$ 内有聚点，并且对每个 $z\in E$，数列
$(f_n(z))$ 收敛，则存在全纯函数 $f$，使 $f_n\to f$ 在 $\Omega$ 上局部
一致。
\end{theorem}

\begin{proof}
由 Montel 定理，$(f_n)$ 至少有一个局部一致收敛子列，记其极限为 $f$。
若另一子列有局部一致极限 $g$，则对每个 $z\in E$，两者都必须等于原数列
$f_n(z)$ 的同一极限。因此 $f=g$ 于 $E$，而 $E$ 在 $\Omega$ 内有聚点；
恒等定理给 $f\equiv g$。

现在反设原序列不局部一致收敛到 $f$。则存在紧集 $K\Subset\Omega$、
$\varepsilon>0$ 和一个子列 $(f_{n_j})$，使
\[
 \sup_K|f_{n_j}-f|\ge\varepsilon\qquad(j\ge1).
\]
该子列仍局部有界，Montel 定理给出进一步局部一致收敛子列。按上一段，它的
极限只能是 $f$，这与所写下界矛盾。因此整个序列局部一致收敛到 $f$。
\end{proof}

局部有界性不可省略。在单位圆盘中令
\[
 p_n(z)=n!\, z\prod_{k=1}^{n}\left(z-\frac1k\right).
\]
对 $E=\{0\}\cup\{1/k:k\ge2\}\subset\mathbb D$ 中每一点，$p_n(z)$ 最终
为零，故逐点收敛；$E$ 的聚点 $0$ 位于圆盘内。但是在圆盘内的点
$z=-1/2$，
\[
 |p_n(-1/2)|
 =\frac{n!(n+1)(n+2)}{2^{n+2}}\longrightarrow\infty.
\]
所以该列在 $-1/2$ 的邻域不可能局部有界，也不可能局部一致收敛。

集合 $E$ 的聚点必须位于区域内部。令
\[
 g(z)=\sin\frac{1}{1-z},\qquad
 E=\left\{1-\frac{1}{k\pi}:k\ge1\right\}\subset\mathbb D,
\]
并令 $f_{2m}=0$、$f_{2m+1}=g$。这个函数列只含两个全纯函数，当然局部有界；
它在 $E$ 上恒为零，故逐点收敛。然而 $E$ 在 $\mathbb D$ 内没有聚点，在
$g(z)\ne0$ 的点处数列在 $0$ 与 $g(z)$ 之间交替，因而不收敛。边界点 $1$
不能供恒等定理使用。

\section{全纯函数族的紧性判据与应用}

Montel 定理常通过容易核验的充分条件使用。第一种情形是值域有界：若所有
$f\in\mathcal F$ 都把 $\Omega$ 映入同一个有界集合，则 $\mathcal F$ 正规。
第二种情形是导数受控并固定一个基点。

\begin{proposition}[导数判据]\label{prop:normal-derivative-criterion}
设 $a\in\Omega$。若 $\{f(a):f\in\mathcal F\}$ 有界，并且导函数族
$\{f':f\in\mathcal F\}$ 局部有界，则 $\mathcal F$ 正规。
\end{proposition}

\begin{proof}
对非空紧集 $K\Subset\Omega$，取有限个圆盘 $B(z_j,r_j)$ 覆盖 $K$，
且各闭圆盘包含于 $\Omega$。用域内折线 $\alpha_j$ 把 $a$ 连到 $z_j$，
令 $L$ 为这些折线的像及闭圆盘的有限并。则 $L$ 紧且包含于 $\Omega$，
导函数族在其上有统一界 $M_L$。
对每个 $z\in K$，选含有它的圆盘，先沿 $\alpha_j$ 再沿直线段走到 $z$。
这条路径全在 $L$ 内，长度不超过
\[
 C_L=\max_j\bigl(\ell(\alpha_j)+r_j\bigr).
\]
沿路径积分，得到
\[
 |f(z)|\le |f(a)|+C_LM_L\qquad(z\in K).
\]
故 $\mathcal F$ 局部有界，再用 Montel 定理。
\end{proof}

正规族还提供极值对象。设 $\mathcal F$ 非空、局部有界，并且在局部一致收敛下
封闭；设 $J:\mathcal F\to\R$ 连续。先说明 $J$ 有界。否则可选
$|J(f_n)|\to\infty$，而 Montel 定理及封闭性给出局部一致收敛到
$f\in\mathcal F$ 的子列，连续性又使该子列的泛函值趋于有限数 $J(f)$，矛盾。
于是上确界有限。取泛函值趋于上确界的序列，再抽取收敛子列，连续性就说明
极限取得上确界；下确界同理。

单叶函数族一般并不封闭，不能直接套用这段论证：它的极限可能成为常数。
若极值序列在固定点的导数模有正下界，导数收敛便排除常值极限，再由单叶极限
定理使极限留在所需函数类中。下一章证明 Riemann 映射定理时会具体核验这一点。

\section*{习题}
\addcontentsline{toc}{section}{习题}

\subsection*{A组}
\begin{enumerate}[label=\textbf{\thechapter.\arabic*.}]
\exerciseitem{normal-local-uniform-integrals}设连续函数 $f_n$ 在 $\Omega$ 上局部一致收敛到 $f$。证明沿任意固定的域内分段 $C^1$ 曲线，积分也收敛。
\exerciseitem{normal-derivative-convergence}设全纯函数 $f_n$ 在 $\Omega$ 上局部一致收敛到 $f$。由 Cauchy 积分公式推导 $f_n'\to f'$ 在紧集上的显式误差估计。
\exerciseitem{normal-power-sequence}证明 $z^n$ 在单位圆盘上局部一致趋于零，但在整个圆盘上不一致收敛。
\exerciseitem{normal-compact-open-metric}验证正文中的度量 $d$ 恰好诱导局部一致收敛，并说明更换紧集穷竭不改变收敛概念。
\end{enumerate}

\subsection*{B组}
\begin{enumerate}[label=\textbf{\thechapter.\arabic*.},resume]
\exerciseitem{normal-cauchy-equicontinuity}补全 Montel 定理中由 Cauchy 估计推出等度连续的有限覆盖论证。
\exerciseitem{normal-disk-maps}设 $\mathcal F\subset\mathcal O(\Omega,\mathbb D)$。证明它正规，并刻画局部一致极限可能落在单位圆周上的情形。
\exerciseitem{normal-basepoint-derivatives}证明命题\ref{prop:normal-derivative-criterion}在 $\Omega$ 非凸时仍成立。
\exerciseitem{normal-vitali-dense-set}设局部有界全纯函数列在 $\Omega$ 的一个稠密子集上逐点收敛。证明它局部一致收敛。
\exerciseitem{normal-vitali-necessity}详细验证正文的多项式列 $p_n$ 说明 Vitali--Porter 定理不能删去局部有界性。
\end{enumerate}

\subsection*{C组}
\begin{enumerate}[label=\textbf{\thechapter.\arabic*.},resume]
\exerciseitem{normal-finite-normal-converse}不用命题\ref{prop:normal-iff-locally-bounded}，直接从正规性的序列定义证明紧集上一致有界。
\exerciseitem{normal-holomorphic-series}设各 $f_n$ 在区域 $\Omega$ 上全纯，$\sum f_n$ 的部分和在 $\Omega$ 上局部有界，并在具有内部聚点的集合上逐点收敛。证明该级数局部一致收敛，且可逐项求导。
\exerciseitem{normal-coefficient-compactness}设 $f_n:\mathbb D\to\mathbb D$，$f_n(z)=\sum_{k\ge0}a_{n,k}z^k$。证明可取子列使每个系数 $a_{n,k}$ 收敛，并说明所得系数确实来自一个全纯极限。
\optionalexerciseitem{normal-univalent-closure}设单叶函数 $f_n:\Omega\to\mathbb D$ 局部一致收敛且 $|f_n'(a)|\ge c>0$。证明极限仍单叶，并解释导数下界的作用。
\end{enumerate}
